\begin{definition}[Topological rings]
\label{definition-topological-ring}
\begin{reference}
\cite[Chapter 0, Sections 7.1 and 7.2]{EGA1}
\end{reference}
Let $R$ be a ring and let $M$ be an $R$-module.
\begin{enumerate}
\item We say $R$ is a {\it topological ring} if $R$ is endowed with a topology
such that both addition and multiplication are continuous as maps
$R \times R \to R$ where $R \times R$ has the product topology.
In this case we say $M$ is a {\it topological module} if $M$ is endowed
with a topology such that addition $M \times M \to M$ and
scalar multiplication $R \times M \to M$ are continuous.
\item A {\it homomorphism of topological modules} is just a continuous
$R$-module map. A {\it homomorphism of topological rings} is a
ring homomorphism which is continuous for the given topologies.
\item We say $M$ is {\it linearly topologized} if $0$ has a fundamental
system of neighbourhoods consisting of submodules. We say $R$ is
{\it linearly topologized} if $0$ has a fundamental system of neighbourhoods
consisting of ideals.
\item If $R$ is linearly topologized, we say that $I \subset R$ is an
{\it ideal of definition} if $I$ is open and if every neighbourhood
of $0$ contains $I^n$ for some $n$.
\item If $R$ is linearly topologized, we say that $R$ is {\it pre-admissible}
if $R$ has an ideal of definition.
\item If $R$ is linearly topologized, we say that $R$ is {\it admissible} if
it is pre-admissible and
complete\footnote{By our conventions this includes separated.}.
\item If $R$ is linearly topologized, we say that $R$ is {\it pre-adic} if
there exists an ideal of definition $I$ such that $\{I^n\}_{n \geq 0}$
forms a fundamental system of neighbourhoods of $0$.
\item If $R$ is linearly topologized, we say that $R$ is {\it adic} if
$R$ is pre-adic and complete.
\end{enumerate}
Note that a (pre)adic topological ring is the same thing as a (pre)admissible
topological ring which has an ideal of definition $I$ such that $I^n$ is
open for all $n \geq 1$.
\end{definition}